\documentclass[11pt]{article}
\usepackage{icmlko}
\begin{document}

\title{판 하나만 속는다 --- 여섯 자 중 다섯이 반대로 말한다}
\author{월드모델 실험실}
\date{노트 474 · 2026년 8월 3일}
\maketitle

\begin{abstract}
노트 473 이 물음 하나를 남겼다 --- \emph{깊이 $12$ 는 판과 무리 안 판을
\textbf{둘 다} 잃는데 \texttt{it}$\cdot$lr 은 무리 안에서만 잃는다. 왜
깊이만 다른가.} 지도가 비어 있어서 못 답했다: 깊이는 자 넷이 다 있는데
\texttt{it}$\cdot$lr 은 집 밖 둘이 비어 있었다.

채웠다.

\begin{center}
\begin{tabular}{lrrrrrr}
\hline
후보 $-$ 챔피언 & 판 & 무리 안 & KR 전체 & KR 갈래안 & 앱 전체 & 앱 갈래안 \\
\hline
깊이 $12$ (채택) & $-0.0016$ & $-0.0023$ & $\mathbf{+0.0064}$ & $-0.0090$ & $\mathbf{+0.0061}$ & $+0.0045$ \\
\textbf{\texttt{it} $400$} & $+0.0003$ & $\mathbf{-0.0024}$ & $\mathbf{-0.0085}$ & $\mathbf{-0.0039}$ & $-0.0003$ & $-0.0019$ \\
lr $0.16$ & $+0.0006$ & $-0.0016$ & $-0.0042$ & $-0.0042$ & $+0.0013$ & $-0.0004$ \\
\hline
\end{tabular}
\end{center}

\textbf{\texttt{it} $400$ 은 여섯 중 다섯에서 음수}이고 그중 셋은 구간이
$0$ 을 안 문다. \textbf{통과시키는 자는 판 하나뿐이다.}
\end{abstract}

\begin{figure}[t]
\centering
\includegraphics[width=\textwidth]{figs/f1.pdf}
\caption{초록이 벌고 빨강이 잃는다. 테두리는 구간이 $0$ 을 안 무는 칸.
\texttt{it} $400$ 줄에서 초록에 가까운 칸은 맨 왼쪽 하나다.}
\end{figure}

\section{그래서 깊이는 왜 다른가}

\textbf{둘은 아예 다른 종류의 거래를 한다.}

\begin{quote}
\textbf{깊이} --- 집안(판 · 무리 안 판)을 조금 내주고 \textbf{집 밖
전체}를 산다. 노트 457 이 그것을 알고 채택했다. \emph{바꾸는 거래}다. \\
\textbf{\texttt{it}} --- 어느 자에서도 안 산다. 다섯에서 잃고 판에서만
비긴다. \emph{바꾸는 게 아니라 능력을 판이 보상하는 쪽(무리 사이 순위)
으로 옮길 뿐}이다.
\end{quote}

노트 461 이 이 둘을 \textbf{배우는 용량}과 \textbf{외우는 용량}으로
불렀다. 그때 근거는 학습 적합도였다(\texttt{it} 은 $+0.124$ 사고
깊이는 $+0.024$). 이제 \textbf{여섯 자에서 같은 결론이 나온다} --- 외우는
용량이 사는 것은 \emph{판에만 보이는 것}이다.

\section{판이 유일한 자였다면 넣었을 것이다}

이것이 이 노트의 실무적 요점이다. \texttt{it} $400$ 은 판에서
$+0.0003$ $[-0.0005,+0.0011]$ 로 비긴다. \textbf{노트 378 의 조건 ①은
비기는 변경을 안 막는다}(부호 $\ge 9/12$ 도 $\le 3/12$ 도 아니면 다음
조건으로 넘어가고, 집 밖이 안 움직이면 ②③④도 침묵한다). 판만 봤다면
``무해하니 넣자''가 됐을 것이다.

\textbf{제일 세게 잡는 자는 무리 안 판이 아니라 KR 수집 전체다}
($-0.0085$). 노트 473 이 ``용량 변경은 무리 안 판을 같이 재라''고 적었는데
\textbf{그 조항을 넓혀야 한다.}

\begin{quote}
\textbf{조항(노트 473 을 넓힘)} --- 용량을 늘리는 변경이 판에서
\emph{비기면}, \textbf{다른 자 최소 둘}(무리 안 판 · 집 밖 짝)로 확인하고
그중 하나라도 구간이 $0$ 밖에서 음수면 안 넣는다. \emph{판이 침묵하는
것은 무해하다는 뜻이 아니다.}
\end{quote}

\section{lr $0.16$ 은 이야기가 없다}

정직하게 적는다. lr $0.16$ 은 여섯 자 \textbf{전부} 구간이 $0$ 을 문다
--- 판의 $+0.0006$ 도, 무리 안의 $-0.0016$ 도 비긴 것이다. 노트 459 가
lr $0.16$ 을 ``판은 평평한데 KR 을 잃는다''고 읽었는데, 씨앗 여덟을
짝지어 보니 \textbf{KR 손해도 구간이 $0$ 을 문다}. \emph{lr $0.16$ 은
그냥 아무 일도 안 하는 값이다} --- 기각은 유지되지만 이유가 ``나쁘다''가
아니라 ``같다''로 바뀐다.

\section*{남은 것}

\begin{itemize}
\item \textbf{판이 비기는 다른 변경들} --- K $64$ 도 여섯 자에서 비겼다
(노트 456 · 473). 진짜 무해한 것과 \texttt{it} $400$ 처럼 판만 속는 것을
가르려면 \emph{다른 자를 항상 재는} 수밖에 없다. 값이 싼 자가 필요하다.
\item \textbf{무리 안 판을 못 재는 네 도메인}(도서 · 시장팝업 · 아이돌 ·
팝업)에서는 이 조항이 반만 걸린다.
\item \textbf{배포 범위} --- 여전히 사람 몫.
\end{itemize}

\end{document}
